\section{Conclusion and Outlook}
\label{sec:conclusion}

\subsection{The single object}

Volume~I's bar complex $\barB^{\mathrm{ch}}(\cA)$ and Volume~II's
Swiss-cheese operad $\SCchtop$ are two views of the same structure.
The MC element $\alpha_T$ in the Swiss-cheese convolution algebra
encodes every construction in this volume as a projection: the
$\Ainf$ operations as its closed face, the PVA bracket as its
cohomological shadow, the spectral $R$-matrix as its mixed-color
component, the bulk-boundary-line triangle as the two-color
decomposition of its MC equation, and the genus tower as its
$\hbar$-expansion.  Each of the following results is a coordinate
expression for this single object.

\begin{enumerate}[label=(\Roman*)]
\item \textbf{Axiomatic Definition} (Section~\ref{sec:foundations}): We constructed the HT operad $\OHT=C_\ast\!\bigl(W(\SCchtop)\bigr)$ and defined $A_\infty$ chiral algebras as algebras over $\OHT$, with explicit sesquilinearity axioms relating $\partial$-equivariance to the $\lambda$-calculus.

\item \textbf{Chain-Level Structure} (Section \ref{sec:FM_calculus}): We proved that bulk local operators form an $A_\infty$ chiral algebra at the chain level, with $A_\infty$ identities arising from Stokes' theorem on FM compactifications and Arnold cancellations at codimension-2 corners.

\item \textbf{PVA Descent} (Section \ref{sec:PVA_descent}): We established that cohomology $H^\bullet(\mathcal{A}, Q)$ is a Poisson vertex algebra by verifying all axioms from Khan--Zeng \cite{KZ25} and proving higher operations $m_{k \geq 3}$ vanish in cohomology via explicit homotopies.

\item \textbf{Worked Examples} (Section \ref{sec:examples_complete}): We computed $A_\infty$ operations explicitly for free chiral multiplets, Landau--Ginzburg models (resolving the $m_{k \geq 4}$ truncation contradiction), and abelian Chern--Simons theory.

\item \textbf{Chiral Hochschild Cohomology} (Section \ref{sec:chiral_hochschild}): We connected bulk observables to chiral Hochschild cochains via dimensional descent from 2d Poisson sigma models and factorization homology on the Ran space.

\item \textbf{Bar--Cobar Duality} (Section \ref{sec:bar_cobar}): We extended Gui--Li--Zeng quadratic duality to the $A_\infty$ setting, proving $d_{\overline{B}}^2 = 0$ using FM--Stokes--Arnold calculus and establishing compatibility with GLZ.

\item \textbf{Line Operators and Yangians} (Section \ref{sec:line_operators}): We proved line operators are modules for the Koszul dual $\mathcal{A}^!$, which carries a dg-shifted Yangian structure with spectral $R$-matrix $r(z)$ satisfying the $A_\infty$ Yang--Baxter equation.
\end{enumerate}

\subsection{Concrete Consequences and Open Problems}

The constructions of this volume, combined with Volume~I's five main theorems, enable the following specific results and identify the following open problems.

\begin{enumerate}[label=(\arabic*)]
\item \textbf{Jones polynomial from the spectral $R$-matrix.}
  The spectral $R$-matrix of Section~\ref{sec:spectral_braiding},
  specialized to Chern--Simons theory with gauge group
  $\mathrm{SU}(2)$ at level~$k$, should recover the colored Jones
  polynomial at $q = e^{2\pi i/(k+2)}$ via the Kontsevich integral.
  The bar-cobar framework provides the chain-level model: the
  $\Ainf$ operations $\{m_k\}$ on $\C \times \R$ produce the
  iterated integrals entering the Kontsevich integral, and the
  $\Eone$-coalgebra structure along~$\R$ implements the group-like
  expansion.  Verifying this identification at all levels is
  equivalent to the genus-$0$ MC5 bridge
  (Theorem~\ref{thm:mc5-genus-zero-bridge}) for
  $\widehat{\mathfrak{su}(2)}_k$.

\item \textbf{MC5 at all genera: resolved.}
  At genus~$0$, the Feynman-diagrammatic $m_k$ coincide with the bar
  differential of Volume~I
  (Theorem~\ref{thm:mc5-genus-zero-bridge}).  At genus~$g \geq 1$,
  the bar differential acquires curvature
  $d^2 = \kappa(\cA) \cdot \omega_g$, and the non-vanishing of higher
  $A_\infty$ operations at the chain level is this curved bar structure.
  The extension to all genera is now established in Volume~I via
  inductive genus determination
  (Volume~I, Theorem~\ref*{thm:inductive-genus-determination}),
  $2$D convergence
  (Volume~I, Proposition~\ref*{prop:2d-convergence}),
  the analytic-algebraic comparison
  (Volume~I, Theorem~\ref*{thm:analytic-algebraic-comparison}),
  and general HS-sewing
  (Volume~I, Theorem~\ref*{thm:general-hs-sewing}).

\item \textbf{Drinfeld--Sokolov reduction commutes with the
  Swiss-cheese structure.}
  The $\mathcal{W}$-algebra computations of
  Section~\ref{sec:W-algebras} demonstrate that the bar-cobar
  adjunction is compatible with Drinfeld--Sokolov reduction for
  $\widehat{\mathfrak{sl}}_N$.  This compatibility, once
  established at the $\Ainf$ level, would allow the transport of
  Volume~I's Theorem~B (bar-cobar inversion) through the BRST
  functor, yielding quasi-isomorphisms
  $\Omegach(\barB(\mathcal{W}_N)) \xrightarrow{\sim} \mathcal{W}_N$
  for all principal $\mathcal{W}$-algebras.

\item \textbf{Higher-genus $A_\infty$ with spectral parameters.}
  This volume works on $\C \times \R$ (genus~$0$ in~$\C$).
  The monograph's genus tower
  ($\kappa(\cA)$, $\Theta_\cA$, complementarity)
  should lift to an $\Ainf$ structure with spectral parameters
  on the Ran space of a higher-genus curve.  The curved
  Swiss-cheese algebras of
  Theorem~\ref{thm:curved-swiss-cheese} provide the target:
  the curvature $\kappa(\cA) \cdot \omega_g$ is the obstruction
  to extending the genus-$0$ PVA to higher genus, and
  Volume~I's complementarity theorem (Theorem~C) constrains
  the deformation space.

\item \textbf{Spectral braiding beyond the evaluation locus.}
  The spectral $R$-matrix $R(z)$ of
  Section~\ref{sec:spectral_braiding} is constructed at the
  evaluation locus.
  Promoting $R(z)$ to the full factorization category would advance
  the monograph's MC3 programme (proved in type~$A$; open beyond)
  (Drinfeld--Kohno extension beyond the
  evaluation-generated core).  The chiral Hochschild complex of
  Section~\ref{sec:chiral_hochschild} provides the deformation-theoretic
  framework: obstructions to extending the braiding live in
  $\mathrm{ChirHoch}^3(\cA)$.

\item \textbf{The screening domination problem.}
  The mass-gap and confinement results of
  Section~\ref*{sec:ym-synthesis} are conditional on screening
  positivity.  Resolving this for $\mathfrak{g} = \mathfrak{su}(3)$
  would give the first rigorous chain-level model for
  color confinement via $\Ainf$ chiral algebras.  The relevant
  spectral data is the $R$-matrix at the critical level
  $k = -h^\vee$, where the Feigin--Frenkel center
  provides an infinite supply of commuting integrals of motion.

\item \textbf{Complete strictification and the Kac--Moody frontier.}
  The spectral Drinfeld class controlling the passage from dg-shifted
  Yangians to strict factorization quantum groups vanishes at every
  filtration for every simple Lie algebra
  (Theorem~\ref{thm:complete-strictification}).  The mechanism is root
  multiplicity one: $\dim\fg_\alpha=1$ forces every root-sector
  multilinear Lie space to be one-dimensional, and the BCH rigidity
  (coefficient $1/n$ at filtration~$n$, via the Dynkin projection and
  the beta integral $\int_0^1(1-t)^{n-1}\,dt=1/n$) leaves no free
  cohomological parameter.  The genuine frontier is Kac--Moody
  algebras with root multiplicities $>1$, where higher-dimensional root
  spaces could carry independent obstruction classes.

\item \textbf{Chiral Koszulness from physics: the loop-order criterion.}
  The loop-order spectral sequence on the bar complex
  (Theorem~\ref{thm:loop-spectral-sequence}) gives a width bound:
  $H^{p,q}(\overline{B}^{\mathrm{ch}}(\cA))=0$ for
  $|q|>\mathfrak{L}(\cA)$ (Theorem~\ref{thm:width-bound}).  One-loop
  exact BV quantization implies chiral Koszulness of the pre-reduction
  boundary algebra (Theorem~\ref{thm:one-loop-koszul}), and
  Drinfeld--Sokolov reduction generically destroys Koszulness by
  transferring infinitely many higher operations from the BRST homotopy
  (Theorem~\ref{thm:ds-koszul-obstruction}).  The criterion is sharp
  for the affine lineage; extending it to non-affine HT theories requires
  the jet-order programme of
  Chapter~\ref{chap:affine-half-space-bv}.

\item \textbf{The nonlinear half-space programme: jet-order induction.}
  The reflected weight identity decomposes by jet order
  (Theorem~\ref{thm:jet-inductive-rwi}): the affine case (jet
  order~$1$) is proved
  (Theorem~\ref{thm:affine-half-space-bv}), the quadratic case (jet
  order~$2$) is conjectured with three independent lines of evidence
  (Conjecture~\ref{conj:quadratic-rwi}), and the cubic case (jet
  order~$3$) is the first where genuine cancellation between graph
  types is required (Conjecture~\ref{conj:cubic-rwi}).  The general
  conjecture reduces to a countable tower of identities, each
  involving finitely many reflected graph types.
\end{enumerate}

\subsection{The cross-volume bridge}

The five theorems of Volume~I and the Swiss-cheese structure of
Volume~II are projections of the same MC element:
Theorem~A specializes to $\SCchtop$-bar-cobar when the curve is
$\C$ and the topological direction is~$\R$;
the spectral $R$-matrix recovers the DK-0 shadow of the
Drinfeld--Kohno ladder via the Laplace transform of the
$\lambda$-bracket;
chiral Hochschild cohomology is the physical origin of Theorem~H;
the curved genus tower is Theorem~D in its physical incarnation.
The complete strictification theorem
(Theorem~\ref{thm:complete-strictification}) adds a sixth bridge:
the spectral Drinfeld class, which a priori could obstruct the
passage from quasi- to strict factorization quantum groups,
vanishes identically---the rigidity is not imposed but inherited
from the root-multiplicity-one structure of simple Lie algebras.
This is the ordered/open-sector analogue of homotopy-Koszulity for
the closed sector: both are structural consequences of the underlying
Lie theory, not of the specific HT dynamics.

\begin{remark}[The three colors and their purity]
\label{rem:three-colors-purity}
The Swiss-cheese operad $\SCchtop$ has three structural components:
the closed color $\FM_k(\C)$ (holomorphic), the open color $E_1$
(topological), and the bulk-boundary interaction (the mixed
composition maps).  The three frontier results of this volume are the
obstruction theories for these three components:
\begin{center}
\begin{tabular}{lll}
\textbf{Color} & \textbf{Obstruction} & \textbf{Purity input} \\
\hline
Closed ($\FM_k(\C)$) & bar cohomology width & one-loop exactness \\
Open ($E_1$) & spectral Drinfeld class & root multiplicity one \\
Mixed (bulk-boundary) & reflected weight identity & jet-order convergence
\end{tabular}
\end{center}
The three inputs on the right are different manifestations of
one-dimensionality: one-loop exactness means the effective BV action
is determined by a single loop order; root multiplicity one means
$\dim\fg_\alpha=1$; jet-order convergence means the reflected
identity is controlled one jet order at a time.

The asymmetry between the three colors is the deep content.
The \emph{open-color} obstruction (spectral Drinfeld class) vanishes
universally for all simple Lie algebras, because root multiplicity one
is a property of the root system---it depends on $\fg$ alone, not on
the specific chiral algebra or HT theory.
The \emph{closed-color} obstruction (Koszulness) can fail, because
bar cohomology depends on the specific algebra $\cA$: Kac--Moody is
Koszul (one-loop exact), but Virasoro is not (DS reduction introduces
higher operations).  The \emph{mixed-color} obstruction (reflected
weight identity) is proved for affine PVAs and open beyond.

This is the Chriss--Ginzburg content of the Swiss-cheese structure.
The Steinberg variety presents the Hecke algebra because the geometry
of $G/B$ is rigid enough (purity of IC complexes) that the convolution
algebra is forced.  The bar complex presents the Swiss-cheese algebra
because the Lie theory of $\fg$ is rigid enough (root multiplicity one)
that the open color is forced, while the closed color retains algebraic
freedom.  Strictification is geometric; Koszulness is algebraic.
The bar complex carries both.
\end{remark}

\begin{remark}[The graph-completed frontier]
\label{rem:graph-completed-frontier}
Volume~I's nonlinear modular shadows appendix now extends the proved
shadow tower into a coherent frontier programme.  Three objects
organize the forward attack:
\begin{enumerate}[label=\textup{(\roman*)}]
\item the \emph{modular chiral Feynman
  transform}~$\mathfrak{F}_{\mathrm{mod}}(\cA)$, a graph-completed
  dg object carrying the Maurer--Cartan master
  element~$\boldsymbol\Theta_\cA$
  (Theorem~\ref*{thm:mc2-bar-intrinsic}) whose projections recover
  the entire modular characteristic hierarchy;
\item the \emph{polarized modular graph algebra}
  $\mathcal{G}^{\pm}_{\mathrm{mod}}(\cA)$ with bipartite
  same-colour vanishing, whose \emph{resonance scattering diagram}
  $\mathfrak{Scat}(\cA)$ governs wall-crossing of nonlinear
  modular data;
\item the \emph{complementarity Legendrian}
  $\Lambda_\cA = j^1 S_\cA$, which contactifies the
  complementarity potential and gives geometric meaning to the
  resonance tower as front-singularity discriminants.
\end{enumerate}
For this volume, the immediate consequence is that the modular
wavefunction $\Psi_\cA$ of the quantization programme is controlled by
the graph-completed master element: its semiclassical expansion
produces the genus tower, and its local singularity type is determined
by the nonlinear shadow archetype (Gaussian, Airy, Pearcey, or mixed
catastrophe).  The curved Swiss-cheese structure of
Theorem~\ref{thm:curved-swiss-cheese} provides the physical
incarnation of the edge-insertion differential~$d_{\mathrm{edge}}$.
\end{remark}

\subsection{The modular holography programme}
\label{subsec:modular-holography-programme}

The constructions of this volume, together with Volume~I's modular
engine, suggest that holomorphic-topological holography is not merely a
genus-zero boundary VOA phenomenon but a \emph{modular} phenomenon: the
boundary chiral algebra is the shadow of a modular Maurer--Cartan
problem, and the holographic dictionary lifts from the disk to all genera
via modular Maurer--Cartan theory.

The compact slogan is:
\[
  \text{stable-graph summation}
  \;+\;
  \text{logarithmic configuration geometry}
  \;+\;
  \underbrace{\text{strict factorization quantum group}}_{\text{Thm~\ref{thm:complete-strictification}}}
  \;=\;
  \text{modular holography}.
\]

Three guiding questions structure the programme.

\begin{enumerate}[label=\textup{(\roman*)}]
\item \textbf{Existence.}
  \ClaimStatusConjectured{}
  Can the genus-zero package
  $\mathcal{G}_0(\mathcal{T};\mathcal{B})$---assembled from CDG
  boundary chiral algebras, Zeng's KK reduction, GKW higher operations,
  and DNP dg-shifted Yangians---be completed to a full modular homotopy
  type via bar-cobar or deformation-theoretic completion over the modular
  operad $C_\bullet(\overline{\mathcal{M}}_{g,n})$?

\item \textbf{Geometry.}
  \ClaimStatusConjectured{}
  Can the gluing operations be realized by explicit integrals over
  logarithmic Fulton--MacPherson spaces and their degeneration
  correspondences (cf.\ Mok)?

\item \textbf{Computation.}
  \ClaimStatusConjectured{}
  Can one compute the first genuinely modular coefficients of the
  genus-refined Yangian kernel
  $R_{\mathcal{T}}^{\mathrm{mod}}(z;\hbar)$ in the three basic
  examples: free $\beta\gamma$, Kodaira--Spencer holography, and the
  $\mathcal{W}_3$ Poisson sigma model?
\end{enumerate}

\noindent
The free $\beta\gamma$ system is the strict free test case (genus
dependence enters only through determinant lines and factorization
homology).  Kodaira--Spencer holography is the holographic open/closed
test case with explicit universal defect algebra.  The $\mathcal{W}_3$
PVA is the first nonlinear higher-spin arena where modular dg-shifted
Yangians genuinely meet stable-curve gluing.

\begin{remark}[A fourth worked example: the M2-brane model]
\label{rem:m2-brane-model}
Costello's M2-brane example provides a fourth arena that exercises the full modular
machine.  Starting from 11d twisted M-theory, one reduces to 5d noncommutative
Chern--Simons theory on $\R \times \C^2$.  The boundary algebra is a \emph{quantum
double-loop algebra}: the unique filtered quantization of
$\mathfrak{gl}_r \otimes \C[u,v]$, where $(u,v)$ are the coordinates on~$\C^2$.
Its Koszul dual
\[
\mathcal{A}_{M2}^!
\;=\;
H^\bullet\bigl(\bar{B}(\mathcal{A}_{M2,\infty})\bigr)^{\vee}
\]
is the holographic dual obtained by Verdier duality on bar cohomology.
The matching between bulk and boundary is explicit:
\begin{itemize}
\item sector $\longleftrightarrow$ adjoint label,
\item level $\longleftrightarrow$ double-loop bidegree in $(u,v)$,
\item $N$ (rank) $\longleftrightarrow$ trace-relation quotient,
\item genus $\longleftrightarrow$ loop order in the 5d theory.
\end{itemize}
The full modular homotopy type assembles as
\begin{equation}
\label{eq:m2-modular-homotopy-type}
\mathcal{D}_{M2}
\;=\;
d_{A_\infty}^{M2} + d_{\mathrm{sew}}^{M2} + \hbar\,\Delta_{5d} + d_r^{M2},
\end{equation}
with $d_{A_\infty}^{M2}$ the bar differential, $d_{\mathrm{sew}}^{M2}$ the stable-graph
sewing, $\Delta_{5d}$ the BV Laplacian from the 5d bulk, and $d_r^{M2}$ the spectral
twist.  This is the most constructive realization of bar--cobar duality as holographic
duality: the algebraic engine of Volume~I and the 3d HT framework of this volume
converge to a single, computable object.
\end{remark}

The volume's Swiss-cheese recognition theorem
(Theorem~\ref{conj:recognition}), homotopy-Koszulity of
$\SCchtop$ (Theorem~\ref{thm:homotopy-Koszul}), and the Hochschild
bridge (Theorem~\ref{thm:hochschild-bridge-genus0}) provide the
algebraic foundations.  Volume~I's five main theorems---bar-cobar
adjunction (Theorem~A), inversion (Theorem~B), complementarity
(Theorem~C), modular characteristic (Theorem~D), and Hochschild
polynomial (Theorem~H)---provide the engine.  The modular holography
programme is the natural next chapter.

\section*{The Lagrangian principle}

\providecommand{\cL}{\mathcal{L}}
\providecommand{\Mvac}{\cM_{\mathrm{vac}}}
\providecommand{\Mbar}{\overline{\mathcal M}}

The categorical logarithm of Volume~I is a Lagrangian embedding.
The bar construction $\barB^{\mathrm{ch}}(\cA)$ maps a chiral
algebra~$\cA$ into the derived moduli~$\Mvac$ of local bulk vacua,
and the image is a shifted-Lagrangian
$\cL_\cA \hookrightarrow \Mvac$: the $(-2)$-shifted symplectic
structure on $\Mvac$
(Conjecture~\ref{conj:geometric-steinberg}(a)) pulls back to zero
along~$\cL_\cA$, with the vanishing witnessed by the bar
differential.  The Swiss-cheese algebra of Volume~II is the
self-intersection of this Lagrangian:
\[
  \mathfrak{S}_\cA
  \;=\;
  \cL_\cA \times_{\Mvac} \cL_\cA,
\]
the derived fiber product that carries a $(-1)$-shifted symplectic
structure inherited from complementary Lagrangians
(Volume~I, Theorem~C).  Every theorem in both volumes is a shadow of
this geometry.  The bar-cobar adjunction (Theorem~A) is the
Lagrangian neighborhood theorem; Koszul inversion (Theorem~B) is the
formal Darboux isomorphism at the self-intersection; the modular
characteristic $\kappa(\cA) \cdot \omega_g$ (Theorem~D) is the
curvature of the Lagrangian family over~$\Mbar_{g,n}$; the chiral
Hochschild ring (Theorem~H) is the ring of functions on the
self-intersection.  The five theorems are five faces of a single
Lagrangian.

\medskip

\noindent\textbf{The three questions, geometrically.}\enspace
The guiding questions of the modular holography programme
acquire their sharpest form when restated through the Lagrangian lens.

\begin{enumerate}[label=\textup{(\roman*${}'$)}]
\item \textbf{Existence as modular completion.}
  The genus-zero Lagrangian $\cL_\cA \hookrightarrow \Mvac$ is
  constructed from bar-cobar data on
  $\FM_k(\C) \times \mathrm{Conf}_k^{<}(\R)$.  Does this embedding
  extend to a family of Lagrangians
  $\cL_\cA^{(g,n)} \hookrightarrow \Mvac$ parametrized
  by~$\Mbar_{g,n}$, so that the genus-$g$ bar complex with curvature
  $d^2 = \kappa(\cA) \cdot \omega_g$ is the pullback of the
  symplectic structure to the fiber over the generic point of the
  Hodge bundle?  This is the modular completion of the Lagrangian:
  the genus tower is the obstruction sequence for extending an
  isotropic embedding through the boundary strata
  of~$\Mbar_{g,n}$.

\item \textbf{Geometry as Lagrangian descent.}
  The gluing morphisms $\Mbar_{g_1,n_1} \times \Mbar_{g_2,n_2}
  \to \Mbar_{g_1+g_2,n_1+n_2-2}$ are correspondences.  Are these
  correspondences Lagrangian in the appropriate shifted symplectic
  category?  Descent of the bar complex through boundary strata is
  precisely the requirement that the Lagrangian family
  $\cL_\cA^{(g,n)}$ is compatible with Lagrangian composition of
  correspondences---the higher-categorical content of the stable-graph
  sewing~$d_{\mathrm{sew}}$.

\item \textbf{Computation as intersection numbers.}
  Each genus-$g$ coefficient of the modular Yangian kernel
  $R^{\mathrm{mod}}_{\mathcal{T}}(z;\hbar)$ is a derived
  intersection number:
  \[
    R^{(g)}_\cA(z)
    \;=\;
    \chi\!\bigl(
      \cL_\cA^{(g)} \times_{\Mvac} \cL_\cA^{(g)}
    \bigr),
  \]
  the Euler characteristic of the genus-$g$ Lagrangian
  self-intersection.  Computing these numbers in the three basic
  examples ($\beta\gamma$, Kodaira--Spencer, $\mathcal{W}_3$ PSM) is
  the genus-refined analogue of computing Kazhdan--Lusztig
  polynomials from the Steinberg variety.
\end{enumerate}

\medskip

\noindent\textbf{Koszulness as the watershed.}\enspace
In the Lagrangian picture, Koszulness has a precise geometric meaning:
a chiral algebra~$\cA$ is Koszul if and only if $\cL_\cA$ meets
itself cleanly---the derived self-intersection $\mathfrak{S}_\cA$ is
concentrated in the expected degree, with no higher Tor.  This is
transversal Lagrangian intersection: the bar cohomology is generated
in a single internal degree, all higher $\Ainf$ operations are
homotopically trivial, and the genus-$g$ corrections are scalar
multiples of the genus-zero answer.  In the physical language, the
QFT is perturbatively exact; loop corrections renormalize but do not
deform.  Kac--Moody algebras at generic level are Koszul, and their
Lagrangians are transversal---the self-intersection is the
Steinberg variety of the loop group, smooth as a derived scheme.

When Koszulness fails, the Lagrangian self-intersection carries
excess: higher $\Ainf$ operations proliferate because the derived
fiber product has Tor in unexpected degrees.  The curved bar structure
$d^2 = \kappa(\cA) \cdot \omega_1$ is the first-order excess
intersection class.  Curvature and braiding entangle because the
excess Tor classes mix the holomorphic direction (where the
$\lambda$-bracket lives) with the topological direction (where the
$\Eone$-coalgebra lives).  The pole structure of the OPE determines
the intersection type: simple poles give transversal intersection and
Koszul algebras; higher-order poles give excess intersection and
curved $\Ainf$ structures.  Drinfeld--Sokolov reduction manufactures
higher poles from simple ones, which is why it generically destroys
Koszulness---reduction moves the Lagrangian into a position of
excess self-intersection.

\medskip

\noindent\textbf{The Steinberg dictionary as derived symplectic
geometry.}\enspace
Every row of the Steinberg dictionary
(Remark~\ref{rem:steinberg-dictionary}) is a face of the Lagrangian
self-intersection.  The flag variety $G/B$ is the Lagrangian
$\cL_\cA$; the Fulton--MacPherson space $\FM_k(\C)$ is the
Lagrangian in the chiral setting.  The nilpotent cone $\cN$ is the
singular support of the Lagrangian; the collision
divisors~$D_S$ are the singular support in the chiral setting.
The Springer resolution is the bar construction's
compactification; the Steinberg variety is the self-intersection; the
Hecke algebra is the convolution algebra of the self-intersection;
the Swiss-cheese convolution $\Linf$-algebra~$\gSC_T$ is the Hecke
algebra in the chiral setting.  The dictionary is not an analogy.
It is a theorem in derived symplectic geometry: the PTVV
$(-2)$-shifted symplectic structure on~$\Mvac$, the Lagrangian
property of the bar embedding, and the $(-1)$-shifted symplectic
structure on~$\mathfrak{S}_\cA$ from complementary Lagrangians
assemble to produce the convolution algebra by the same formal
mechanism---Borel--Moore homology of the correspondence---that
produces the Hecke algebra from the Steinberg variety.  The algebraic
Steinberg presentation
(Theorem~\ref{thm:steinberg-presentation}) is the shadow on
homology; the geometric Steinberg conjecture
(Conjecture~\ref{conj:geometric-steinberg}) is the full derived
statement.

\medskip

A chiral algebra is a Lagrangian.  The bar complex is its
self-intersection.  The Swiss-cheese operad is the internal structure
of that self-intersection.  Everything else is a shadow.

